\section{Feasibility and distance to target}
\label{sec:feasibility}

The position this paper has been asked to test is that Flux's substrate is already good and the
remaining work is extension rather than invention. That is a testable claim, and this section tests
it rather than repeating it. The result is mixed in a specific and useful way: the cloud layer
largely supports the claim, and four mechanisms do not.

\subsection{Method}

Every mechanism tagged \planned or \vision elsewhere in this paper appears below with what exists
today (cited), what is missing, the hardest unsolved sub-problem, and a risk class. Classes are
assigned from the evidence gathered in Parts II--IV, not from schedule optimism:

\begin{description}
  \item[Extension] composes components that exist, on paths that already run.
  \item[Engineering] a substantial new build, with no unsolved research problem.
  \item[Research] an open problem with no known good answer, or a property that cannot be obtained
    without a new trust assumption.
\end{description}

A note on what an honest assignment costs. If every row came out \emph{Extension}, the exercise would
have failed. Four rows are \emph{Research}, and three of those four are load-bearing for the thesis in
\S\ref{sec:introduction}.

\subsection{The table}

{\footnotesize
\begin{longtable}{@{}p{0.5cm}p{3.0cm}p{3.3cm}p{3.3cm}p{2.0cm}@{}}
\toprule
\# & Mechanism & Exists today & Hardest sub-problem & Class \\
\midrule
\endhead
1 & \PNR settlement and fee pool & \PoN payment rotation, deterministic and \emph{measured} to equalise within one rotation per tier; consensus-enforced coinbase splits already exist & Pool as consensus state with a reorg-exact serialisation (settled: coinbase-committed) and $\drip$ (settled: 86{,}400 blocks) --- \S\ref{sec:pnr} & Engineering \\
2 & Parallel-asset consolidation & Ten deployed contracts; audited Schnorr \mbox{ERC-4337} account on seven EVM chains & Reaching holders of retired assets and handling unclaimed balances --- coordination, not code & Engineering \\
3 & Bounded, publicly verifiable bridge & Audited Schnorr account; Solana multisig live on mainnet & Keeping administrative control decentralized \emph{and} bounded & Engineering \\
4 & Bonded attestation network & Measured boot, dm-verity, LUKS, fail-closed \texttt{fes}; \texttt{purpose=\allowbreak{}mesh} implemented with no consumer & The attestation key is not machine-bound, so a bond would secure a signature anyone with the software can produce & Research $\rightarrow$ Engineering \\
5 & \appspec \specv{9} & 932 files changed, 7{,}763 test cases per push, draining and per-app-CA TLS wired & Migrating 322 deployed applications across six legacy versions without evicting any & Extension \\
6 & Main-chain pruning & \texttt{-prune} fully wired and inherited; \FluxNode registry already in a separate store & Retaining registry and collateral UTXOs across pruning depth & Extension \\
7 & Collateral maturity & Coinbase-maturity precedent; registry enumerable & Defining ``continued participation'' so upgrades stay instant while exits wait & Engineering \\
8 & Post-quantum migration & Nothing: a repo-wide search returns zero hits & Primitive unchosen; residue classes outside collateral. Collateral path designed (\S\ref{sec:pq-frontier}) & Research \\
9 & Moderation quorum & Node owner keys and collateral weights; a \mbox{GitHub-file} blocklist as today's mechanism & Four owner identities (three under linkage) reach the 25\% threshold; no mechanism change fixes concentration & Research \\
10 & Direct per-deployment payment & \emph{Shipped twice}: application registration, and \CumulusVPN's memo-derived entitlement & None; this is the mature path & Extension \\
11 & Deposited credits and renewal & \texttt{appsmonitor} auto-renewal with Foundation-held keys & Authorising a recurring charge without a custodian and without an allowance primitive --- placed outside protocol scope: the protocol offers no auto-renewal, and credits are a third-party custodial service built on \Flux{} \srcintent{08-answers-round-2.md, round 9} & Engineering \\
12 & Agent deployment authority & \FluxOS validates owner signatures over specifications on every node & None; it composes with existing validation & Extension \\
13 & Agent \emph{spending} authority, bounded & \mbox{ERC-4337} session-key substrate on EVM chains; scoped revocable FluxEdge tokens (action-bounded) & Not possible on the \Flux L1 without a new trusted party or new consensus state & Research \\
14 & Agent-native provisioning API & Shipped MCP tools; OpenAPI 3.1.1 FluxEdge API with scoped tokens and auto-revocation & Inherits row 13 for its bounds & Extension \\
15 & Natural-language deployment & \texttt{DeploymentYAML}: retrieval-augmented generation producing Kubernetes YAML; Orbit specification generation & Whether a human can meaningfully audit a generated specification before signing & Engineering \\
16 & Shielded value & Sapling notes, encryption, commitment tree, nullifiers, Groth16 --- all deployed and inherited & None; it runs & Extension \\
17 & Shielded payment for deployments & The primitives above & \emph{Withdrawn}: both pools are to be retired \srcintent{08-answers-round-2.md, round 10} & --- \\
18 & Proving a shielded payment funds a \emph{specific} deployment unlinkably & Note commitments and nullifiers & An application address is itself a workload identifier; at 35.4 renewals/day a payment is a singleton in its block & Research (moot after retirement) \\
19 & Orchard/Halo2 migration & Sapling/Groth16 deployed; no Orchard present & Consensus change plus audit scoping; now a future option rather than a migration & Engineering \\
19b & Shielded-pool retirement & Both pools live, closed to entry, 96{,}479.94\,FLUX & Sunset height, outreach, and removal of the Rust proving stack & Engineering \\
20 & GPU tenancy & Whole-device passthrough via Kubernetes & No MIG, vGPU or time-slicing; requests silently stripped on contention & Engineering \\
21 & \FluxCloud/\FluxEdge convergence & Two runtimes and two products & Organisational as much as technical & Engineering \\
22 & Decentralizing the reachability path & \FDM, \PDNS, gateway and certificates all working & \texttt{api.\allowbreak{}runonflux.\allowbreak{}io} is the sole application-discovery source, with no on-chain or P2P fallback & Engineering \\
\bottomrule
\caption{Distance from the deployed system to the target architecture. Citations for every ``exists
today'' entry are given in the section that describes it.}
\label{tab:feasibility}
\end{longtable}
}

\noindent Tally: \textbf{Extension 6, Engineering 11, Research$\rightarrow$Engineering 1, Research 4.}

\subsection{The four Research items, which are not equally hard}

\paragraph{Post-quantum migration (row 8).}
Hard for everyone, and \Flux{} starts from a \emph{worse} position than Bitcoin on one axis and a
better one on another. Worse: \SI{20.61}{\percent} of supply sits in \FluxNode{} collateral whose
public keys are published by protocol mandate, not by user behaviour
(\Cref{rem:pq-collateral-exposed}) --- and BIP-360's remedy, an output type that hides the key in the
scriptPubKey \citep{bip360}, does not reach it, because the exposure is in the registration
transaction rather than the output script.

Better, and it is the axis that decides the disposition: collateral is a \emph{registered} set with a
mandatory recurring transaction, so \Flux{} has an authenticated channel to exactly the exposed
population. That is what makes eligibility-forced migration possible
(\Cref{rem:pq-disposition}, adopted) and what lets \Flux{} avoid the freeze-or-tolerate-theft dilemma
BIP-361 must resolve for Bitcoin \citep{bip361}. The row stays \textbf{Research} because the
primitive is unchosen and the residue classes are unsolved, not because the collateral path is
undesigned --- that path is specified in \S\ref{sec:pq-frontier}.

\paragraph{Bounded agent spending on the \Flux L1 (row 13).}
Hard for a precise and stateable reason established in \S\ref{sec:agent-identity}: script validation
is a pure function of the spending transaction and the outputs it consumes, so no rule can enforce a
rate limit or a cumulative cap. The deployable answer --- fund a dedicated key, and the budget is the
balance --- is honest and unglamorous. Everything richer requires a co-signer, which is a new trusted
party, or new consensus state.

\paragraph{Unlinkable payment for a specific deployment (row 18).}
The paper's most interesting problem, and the one whose difficulty is easiest to overstate. It must
be classed in three parts and never collapsed: the shielded primitives are inherited and shipped
(Extension); applying them to deployment payment is Engineering, gated on row 17; and proving a
shielded payment funds a \emph{particular} deployment without linking payer to workload is Research.
Crucially, \textbf{it applies only to the optional credit path and to shielded funding proofs}. The
default direct-payment path is unaffected, because the payment \emph{is} the renewal and there is no
standing obligation to prove --- so the agent-native argument of \S\ref{sec:agent-native} survives
even with this problem open. That limitation of blast radius is a real result and not a consolation.

\paragraph{The moderation quorum (row 9).}
Classed Research not because the mechanism is hard to build --- it is straightforward --- but because
\textbf{the property it is meant to provide cannot be obtained by building it}. Four owner identities
--- three under the strongest linkage computable from public data --- clear the 25\% threshold at the
measured distribution. No bond, rate limit or eligibility gate changes
that. The open question is whether a two-chamber rule requiring both a weight threshold and a
distinct-identity count can preserve Sybil resistance while adding censorship resistance, and that is
genuinely open.

\subsection{Where the claim holds, and where it does not}

For the cloud layer the claim holds well. \specv{9}, pruning, deployment authority, direct payment and
the provisioning API are all Extension, and the \specv{9} substrate is demonstrably built and tested.
The proposition that \Flux can extend rather than invent its way to an agent-serving cloud is
\emph{largely correct, for everything except payment bounds and payment privacy}.

Where it does not hold, row 4 is the load-bearing case. The roadmap treats the bonded attestation
network as an extension of an existing attestation service, but the existing service does not provide
the property the bonded design needs to build on: an attestation demonstrates possession of a value
shipped in every copy of the software, not that a given machine booted a given image
(\S\ref{sec:attested-execution}). That is a prerequisite, not an increment, and it should be sequenced
as one.

\subsection{The dependency that reorders the schedule}

One edge in the dependency graph deserves to be lifted out of the table, because it changes what
should ship first.

\PNR pays all nodes from a shared pool, so a node's income does not depend on hosting; marginal
hosting income is exactly zero and free-riding dominates at every parameter value
(\S\ref{sec:pnr}). What compels hosting is therefore the eligibility gate --- and rows 4 and 9
establish that the gate is neither operator-bound nor machine-bound. \textbf{\PNR raises the payoff to
passing that gate without doing the work, by exactly what it pays.} Hardening attestation is thus a
\emph{precondition} for \PNR activation rather than a parallel workstream; sequencing it afterwards
would deploy an incentive to fake eligibility before the mechanism that detects faking exists. This
single edge raises the practical difficulty of row 1 above its own construction class, and
\S\ref{sec:migration} treats it as the schedule's binding constraint.

A second, milder dependency: rows 1, 5, 11, 14 and 21 all wait on \specv{9} fields that currently have
no consumers --- \texttt{paymentMemo}, \texttt{referralCode} and \texttt{machineType}
(\S\ref{sec:appspec-v9}). \specv{9} is therefore the critical path for most of Part IV, and its
enforcement height, since settled at $\hstar = \num{3050000}$ (\S\ref{sec:appspec-v9}), is the
single most schedule-relevant date in this document.

%% Drain this section's float queue before the next begins.
\FloatBarrier
